% !TEX root = ../main.tex
\section{Results and Discussion}
\label{sec:results}

\plantodo{S5}{Present the nominal-flight ConOps (strategic/commercial
planning through day-of-operations, turnaround, flight execution, and
postflight/continuation) as the operational backbone the architecture
models. Summarize; put the full activity breakdown in
Appendix~\ref{app:conops}.}

\plantodo{S8}{Present the stakeholder objective/cost/value ontology:
what "optimal" means per stakeholder (airline, passenger, ATC/ANSP,
airport, flight crew, environmental/societal, military), and the
classification of each objective as hard constraint, optimization
objective, cost/penalty, MOE, or MOP. Highlight conflicting vs. aligned
objectives and objectives with externalized costs.}

\plantodo{S9}{Present the myopic-optimization analysis: the definition of
"local optimum" in the SoS context, candidate SoS measures of value, and
the documented conflicts (decision maker, decision variable, local
objective/constraints, affected stakeholders, externalized costs, SoS
consequence). See the flagged scope note below.}

\textbf{Scope note (flag, not a TODO to resolve silently):} Project\_To-Do
List.md S9 catalogs eight candidate local-vs-SoS conflicts, each with eight
documented attributes (64 cells). A 15--40 page report body cannot carry
all eight in full; plan to present 2--4 illustrative conflicts in the main
body (favoring whichever one becomes the S11--S14 demonstration scenario)
and move the remaining catalog entries to Appendix~\ref{app:myopic}.

\plantodo{S10}{Present the SysML architecture: context diagram, package
structure, stakeholder model, BDDs, IBDs, information-object model,
interface/item-flow model, nominal-flight activity diagram, responsibility
swimlanes, and requirements model with traceability. Embed the most
illustrative diagrams here; move the full diagram set to
Appendix~\ref{app:architecture}.}

\plantodo{S12, S13}{Present the optimization/simulation results: the
baseline scenario, single-stakeholder vs. balanced/SoS optimization
outcomes, and how simulation entities/decisions/objectives/constraints
trace back to SysML blocks, properties, goals, and requirements. Show how
the architecture exposes impacts a standalone optimizer would omit (S13's
key point).}

\textbf{Scope note:} Project\_To-Do List.md S12 as written (Pareto fronts,
full objective-weight sweeps, sensitivity analysis, tipping-point
identification across the whole experiment matrix) reads as a full
optimization research study layered on top of the architecture work. Per
this project's own working assumption ("not an optimization paper";
optimization is a capability inside the architecture, not the whole
subject) and the 15--40 page report constraint, consider scoping S12's
demonstration to one representative scenario with a single weight sweep
and one Pareto-style comparison, rather than the full matrix -- see the
note added to \texttt{knowledge/questions/open-questions.md} under
"Optimization study scope (S11)."

% ---- Placeholders added 2026-10-09 (AI-drafted scaffolding, not report prose).
% No results exist yet. Every cell below is empty on purpose.
\subsection{Experiment Results}
\label{sec:experiment-results}

\begin{table}[h]
\centering
\caption{PLACEHOLDER -- resolution options. A dash means no instruction.
Full version with instruction counts in Appendix~\ref{app:optimization}.}
\label{tab:options}
\begin{tabular}{@{}lllll@{}}
\toprule
Option & AC1 & AC2 & AC3 & AC4 \\
\midrule
O1 & Descend 2,000 ft & -- & -- & -- \\
O2 & -- & Descend 2,000 ft & -- & -- \\
O3 & Climb 2,000 ft & -- & -- & -- \\
O4 & -- & Climb 2,000 ft & -- & -- \\
O5 & Turn behind AC2 & -- & -- & -- \\
O6 & -- & Turn behind AC1 & -- & -- \\
O7 & Reduce speed & -- & -- & -- \\
O8 & Descend 4,000 ft & -- & -- & -- \\
O9 & Descend 2,000 ft & -- & Slow or turn & -- \\
\bottomrule
\end{tabular}
\end{table}

\begin{table}[h]
\centering
\caption{PLACEHOLDER -- outcome of each case at four levels, as a difference
from case A. No simulation has been run; cells are empty.}
\label{tab:outcomes}
\begin{tabular}{@{}p{1.6cm}p{1.6cm}p{2.6cm}p{2.6cm}p{2.8cm}p{2.8cm}@{}}
\toprule
Case & Option chosen & Aircraft (largest cost to any one) & Airline (each
airline's total) & Sector (clearances / effort / risk) & System (total fuel,
total CO2, total cost) \\
\midrule
A &  & reference & reference & reference & reference \\
B &  &  &  &  &  \\
C &  &  &  &  &  \\
\bottomrule
\end{tabular}
\end{table}

\plantodo{S12}{\textbf{Which option each case chose, and why.} Walk
Table~\ref{tab:outcomes} row by row. For case A, name the rule that produced
the choice. For B and C, name the cost that decided it.}

\plantodo{S12}{\textbf{Did the climb rule matter?} With the crossing pair
above FL350, cases A and B could not offer a climb and case C could. Say
whether case C used one, for which aircraft, and what it saved. If C did not
climb either, say so: that is a result too.}

\plantodo{S12}{\textbf{B minus A, and C minus B.} The first is what
automation is worth; the second is what the information exchange is worth.
Compare the size of each with the prior figure of tens of gallons per flight
(\cite{engility2014tasarAlaska}, p.~12) and say whether the difference is
larger than the truth model's uncertainty.}

\plantodo{S9, S12}{\textbf{Does the best option change with the level?} Read
across each row of Table~\ref{tab:outcomes}. Name any option that is better
for one aircraft or one airline and worse for the sector or the system, and
who pays (the largest cost to any one aircraft). This is the sub-system
versus system discussion.}

\plantodo{S12}{\textbf{Airline X's two flights.} Discuss the option that
touches both of one airline's aircraft: what the airline would choose
between its late flight and its early one, and that the controller cannot
see that preference. Connect to the dispatcher priorities (schedule
integrity and connections rated at or above fuel; source of that table still
to confirm).}

\plantodo{S12}{\textbf{Workload is three things.} Discuss clearances issued,
effort to devise, and residual risk separately. Say where a ``more
workload'' option on one of them is ``less'' on another.}

\plantodo{S12}{\textbf{The case where nothing changes.} Report the variation
with similar weights and cost indexes, where case C picks the same option as
case A, to show the comparison is not rigged.}

\plantodo{S13}{\textbf{Trace back to the model.} For each column of
Table~\ref{tab:outcomes}, name the part in the context diagram it belongs to
(aircraft and crew; airline operations center; ATC sector). Say what the
system column maps to, since no single part owns it. Name the exchange that
case C adds (weight and cost index to ATC) and where it sits in the model.}

\plantodo{S12, S15}{\textbf{Limits.} Level crossing conflicts are a minority
(36 of 256, \cite{rantanen2012conflictManeuvers}); one invented sector; a
static snapshot; the climb rule is one threshold for every type; airlines
could misreport what they share (cite the 1 to 3 percent figure once that
source is registered).}
% ---- end placeholders 2026-10-09
